% Related Work (draft)
\label{sec:related}

\subsection{Statistical and Machine-Learning Models for Football}
The canonical generative model for football scores is the Poisson
model of Maher~\citep{maher1982modelling}, refined by
Dixon and Coles~\citep{dixon1997modelling}, who introduce a
low-score correction and time-decayed team parameters. Elo-style
ratings~\citep{elo1978rating} remain a strong and transparent
baseline, and gradient-boosted decision trees have become the
standard tabular approach in sports analytics. A consistent finding
of this literature is that predictive accuracy saturates around
50--55\% for three-way match outcomes, close to the empirical base
rates, and that bookmaker odds are difficult to outperform
\citep{hvattum2010elo, spann2009sports}.

\subsection{Market Efficiency and Closing Odds}
Bookmaker closing odds are widely regarded as near-efficient
aggregators of public information~\citep{levitt2004why, franck2010prediction}.
In particular, the closing line is a stronger predictor than the
opening line, and drift from opening to closing odds contains little
exploitable information once transaction costs are considered. Any
claimed predictive edge must therefore be evaluated against
de-vigged closing odds under an identical protocol; evaluations that
omit this baseline cannot distinguish model skill from market
inefficiency. The present work follows this prescription, and additionally
reports the opening-to-closing drift experiment as a direct test of
market efficiency on the present data.

\subsection{LLM Probabilistic Forecasting}
Large language models have recently been shown to produce calibrated
probabilistic forecasts from textual evidence.
Halawi et al.~\citep{halawi2024approaching} demonstrate that
fine-tuned LLMs approach human expert forecasting on general
prediction markets, and the AIA Forecaster~\citep{alur2025aia}
scales this to a multi-agent retrieval-augmented system.
LLM-as-a-Prophet~\citep{yang2025llm} provides a controlled
``prophet arena'' in which the predictive intelligence of LLMs can be
benchmarked. These systems rely on web retrieval and fine-tuning,
which raises the bar for reproducibility and makes leakage control
difficult to audit. The present setting is deliberately simpler: the LLM
receives only a pre-match feature card, has no retrieval access, and
is evaluated against the market under the same protocol as all other
models. This isolates the contribution of domain-enhanced reasoning
from retrieval and training data effects.

\subsection{Gaps}
Three gaps motivate the design. (i)~Most evaluations lack a
de-vigged market baseline under an identical betting protocol.
(ii)~Leakage controls---pre-match information cutoffs, train-only
statistics, temporal splits---are rarely stated precisely enough to
reproduce. (iii)~Uncertainty and risk are either absent or defined by
hand-crafted scores whose empirical validity is not tested. All
three are addressed with a conservative, fully specified evaluation
and with empirically validated uncertainty signals.
